\documentclass{article}
\usepackage{amsmath, amssymb, amsthm}
\usepackage{hyperref}
\usepackage{geometry}
\geometry{margin=1in}

\title{Erd\H{o}s Problem \#667}
\date{Last updated: 2025-08-31}
\author{Source: erdosproblems.com}

\begin{document}
\maketitle

\noindent\textbf{Status:} OPEN \quad \textbf{No prize}

\noindent\textbf{Tags:} graph theory, ramsey theory

\medskip

\noindent\textbf{Problem Statement:}

Let $p,q\geq 1$ be fixed integers. We define $H(n)=H(N;p,q)$ to be the largest $m$ such that any graph on $n$ vertices where every set of $p$ vertices spans at least $q$ edges must contain a complete graph on $m$ vertices. Is\[c(p,q)=\liminf \frac{\log H(n)}{\log n}\]a strictly increasing function of $q$ for $1\leq q\leq \binom{p-1}{2}+1$?




A problem of Erd\H{o}s, Faudree, Rousseau, and Schelp. When $q=1$ this corresponds exactly to the classical Ramsey problem, and hence for example\[\frac{1}{p-1}\leq c(p,1) \leq \frac{2}{p+1}.\]It is easy to see that if $q=\binom{p-1}{2}+1$ then $c(p,q)=1$. Erd\H{o}s, Faudree, Rousseau, and Schelp have shown that $c(p,\binom{p-1}{2})\leq 1/2$.


\bigskip
\noindent\small{Source: \url{https://www.erdosproblems.com/667}}
\end{document}
